\documentclass[../../main.tex]{subfiles}
\begin{document}

{\bf [解答]}

$k=1,2,3,4$とする．$P_k(a_k,b_k)$とおく．題意から
\begin{align}
  \begin{cases}
    (2a_1,2b_1)=(a_4,b_4+2)  \\
    (2a_2,2b_2)=(a_1,b_1) \\
    (2a_3,2b_3)=(a_2+1,b_2)  \\
    (2a_4,2b_4)=(a_3+2,b_3+2)
  \end{cases}
  \Longleftrightarrow
  \begin{cases}
    2a_1=a_4  \\
    2a_2=a_1 \\
    2a_3=a_2+1  \\
    2a_4=a_3+2
  \end{cases}, \quad
  \begin{cases}
    2b_1=b_4+2  \\
    2b_2=b_1 \\
    2b_3=b_2  \\
    2b_4=b_3+2
  \end{cases}
\end{align}
である．まず$a_k$について
\begin{align}
  &a_1=\dfrac{a_4}{2}=\dfrac{a_3+2}{4}=\dfrac{a_2+5}{8}=\dfrac{a_1+10}{16} \\
  \Longleftrightarrow
  &a_1=\dfrac{2}{3}
\end{align}
となり，順に$(a_2,a_3,a_4)=(1/3,2/3,4/3)$である．次に同様に$b_k$について
\begin{align}
  &b_1=\dfrac{b_4+2}{2}=\dfrac{b_3+6}{4}=\dfrac{b_2+12}{8}=\dfrac{b_1+24}{16}\\
  \Longleftrightarrow
  &b_1=\dfrac{8}{5}
\end{align}
であるから，$(b_2,b_3,b_4)=(4/5,2/5,6/5)$が従う．以上から
\begin{align}
  P_1\left(\dfrac{2}{3},\dfrac{8}{5}\right) , P_2\left(\dfrac{1}{3},\dfrac{4}{5}\right) ,
  P_3\left(\dfrac{2}{3},\dfrac{2}{5}\right) , P_4\left(\dfrac{4}{3},\dfrac{6}{5}\right)
\end{align}
である．$\cdots$(答)

\begin{figure}[htb]
  \centering
  \begin{tikzpicture}[x=3.2cm, y=2.2cm, >=stealth, every node/.style={font=\small}]
    \draw[->, thick, gray!80!black] (-0.2, 0) -- (1.8, 0) node[right] {$x$};
    \draw[->, thick, gray!80!black] (0, -0.2) -- (0, 2.0) node[above] {$y$};
    \node[below left=2pt] at (0, 0) {$O$};

    \coordinate (P1) at ({2.0/3.0}, {8.0/5.0});
    \coordinate (P2) at ({1.0/3.0}, {4.0/5.0});
    \coordinate (P3) at ({2.0/3.0}, {2.0/5.0});
    \coordinate (P4) at ({4.0/3.0}, {6.0/5.0});

    % 領域のハッチング
    \fill[pattern=north east lines, pattern color=blue!40]
    (P1) -- (P2) -- (P3) -- (P4) -- cycle;

    % 四辺形の輪郭
    \draw[line width=1.2pt, blue!85!black] (P1) -- (P2) -- (P3) -- (P4) -- cycle;

    % 対角線 P1 P3 (垂直線)
    \draw[line width=1.0pt, dashed, blue!70!black] (P1) -- (P3);

    % 座標軸への投影点線
    \draw[densely dotted, thick, gray!70!black] (P1) -- (0, {8.0/5.0}) node[left] {$\frac{8}{5}$};
    \draw[densely dotted, thick, gray!70!black] (P2) -- (0, {4.0/5.0}) node[left] {$\frac{4}{5}$};
    \draw[densely dotted, thick, gray!70!black] (P3) -- (0, {2.0/5.0}) node[left] {$\frac{2}{5}$};
    \draw[densely dotted, thick, gray!70!black] (P4) -- (0, {6.0/5.0}) node[left] {$\frac{6}{5}$};

    \draw[densely dotted, thick, gray!70!black] (P2) -- ({1.0/3.0}, 0) node[below] {$\frac{1}{3}$};
    \draw[densely dotted, thick, gray!70!black] (P3) -- ({2.0/3.0}, 0) node[below] {$\frac{2}{3}$};
    \draw[densely dotted, thick, gray!70!black] (P4) -- ({4.0/3.0}, 0) node[below] {$\frac{4}{3}$};

    % 頂点ノード
    \node[circle, fill=blue!85!black, inner sep=1.3pt] at (P1) {};
    \node[above=2pt, blue!85!black] at (P1) {$P_1$};
    \node[circle, fill=blue!85!black, inner sep=1.3pt] at (P2) {};
    \node[above left=1pt, blue!85!black] at (P2) {$P_2$};
    \node[circle, fill=blue!85!black, inner sep=1.3pt] at (P3) {};
    \node[below right=1pt, blue!85!black] at (P3) {$P_3$};
    \node[circle, fill=blue!85!black, inner sep=1.3pt] at (P4) {};
    \node[above right=1pt, blue!85!black] at (P4) {$P_4$};
  \end{tikzpicture}
  \caption{四辺形 $P_1P_2P_3P_4$ の配置と形状}
  \label{fig:quadrilateral_p1234}
\end{figure}

従って四辺形$P_1P_2P_3P_4$の面積は
\begin{align}
  S=\dfrac{8-2}{5}\dfrac{4-1}{3}\dfrac{1}{2}=\dfrac{3}{5}
\end{align}
である．
$\cdots$(答)
\newpage
\end{document}
